% preamble.tex -- house style of Gruppe Jahnel (TU Braunschweig)
\usepackage[T1]{fontenc}
\usepackage{lmodern}
\usepackage[utf8]{inputenc}
\usepackage[british]{babel}
\usepackage{amsmath,amssymb,amsthm,mathtools}
\usepackage[a4paper,margin=1.05in]{geometry}
\usepackage{microtype}
\setlength{\emergencystretch}{1em}
\usepackage{xcolor}
\definecolor{paperlink}{RGB}{0,0,114}
% shared figure colours (roles; every colour is paired with a line style or a label)
\definecolor{figmain}{RGB}{0,62,140}   % raw path, corridor, main objects
\definecolor{figaux}{RGB}{176,84,0}    % escapes, inserted pieces
\definecolor{figgrey}{gray}{0.5}       % guides, closed edges
\usepackage{enumitem}
\setlist[itemize]{leftmargin=2em,itemsep=0.25em,topsep=0.35em}
\setlist[enumerate]{leftmargin=2.25em,itemsep=0.25em,topsep=0.35em}
\usepackage{tikz}
\usetikzlibrary{arrows.meta,positioning,calc,decorations.pathmorphing,fit,backgrounds}
\usepackage{fancyhdr}
\pagestyle{fancy}
\fancyhf{}
\fancyhead[L]{\small Nonuniqueness on nonamenable transitive graphs}
\fancyhead[R]{\small\thepage}
\setlength{\headheight}{14pt}
\usepackage[hyphens]{url}
\usepackage{hyperref}
\hypersetup{colorlinks=true,linktoc=all,linkcolor=paperlink,citecolor=paperlink,
  urlcolor=paperlink,pdfstartview=Fit}
\usepackage[capitalise,nameinlink,noabbrev]{cleveref}
\setlength{\parindent}{1.25em}
\setlength{\parskip}{0pt}
\numberwithin{equation}{section}
\setcounter{tocdepth}{1}

% ---------------------------------------------------------------- theorems
\theoremstyle{plain}
\newtheorem{theorem}{Theorem}[section]
\newtheorem{proposition}[theorem]{Proposition}
\newtheorem{lemma}[theorem]{Lemma}
\newtheorem{corollary}[theorem]{Corollary}
\theoremstyle{definition}
\newtheorem{definition}[theorem]{Definition}
\newtheorem{assumption}[theorem]{Assumption}
\newtheorem{example}[theorem]{Example}
\theoremstyle{remark}
\newtheorem{remark}[theorem]{Remark}

% standing assumptions (A1), (A2): referenced with \cref
\newlist{asslist}{enumerate}{1}
\setlist[asslist]{label=\textup{(A\arabic*)},ref=\textup{(A\arabic*)},leftmargin=3em}
\crefformat{asslisti}{#2#1#3}
\Crefformat{asslisti}{#2#1#3}
\crefrangeformat{asslisti}{#3#1#4--#5#2#6}
\crefmultiformat{asslisti}{#2#1#3}{ and~#2#1#3}{, #2#1#3}{ and~#2#1#3}

% ---------------------------------------------------------------- macros
\newcommand{\PP}{\mathbb P}
\newcommand{\EE}{\mathbb E}
\newcommand{\N}{\mathbb N}
\newcommand{\Z}{\mathbb Z}
\newcommand{\R}{\mathbb R}
\DeclareMathOperator{\Var}{Var}
\DeclareMathOperator{\Cov}{Cov}
% indicator: \mathds 1 (dsfont); bold fallback where dsfont is not installed
\IfFileExists{dsfont.sty}{\usepackage{dsfont}\newcommand{\1}{\mathds{1}}}{\newcommand{\1}{\mathbf{1}}}
\newcommand{\e}{\mathrm{e}}
% \d and \c: upright d and complement sign in math mode, the usual accents in text mode
\let\accentdotbelow\d
\DeclareRobustCommand{\d}{\ifmmode\mathrm{d}\else\expandafter\accentdotbelow\fi}
\newcommand{\dd}{\,\mathrm{d}}
\let\accentcedilla\c
\DeclareRobustCommand{\c}{\ifmmode\mathrm{c}\else\expandafter\accentcedilla\fi}
\DeclarePairedDelimiter{\abs}{\lvert}{\rvert}
\DeclarePairedDelimiter{\norm}{\lVert}{\rVert}
\DeclareMathOperator{\tr}{tr}
\DeclareMathOperator{\Adj}{Adj}
\DeclareMathOperator{\dist}{dist}
\DeclareMathOperator{\diam}{diam}
\DeclareMathOperator{\Aut}{Aut}
\DeclareMathOperator{\Piv}{Piv}
\DeclareMathOperator{\Cay}{Cay}
\newcommand{\given}{\,\lvert\,}
\newcommand{\conn}{\leftrightarrow}
\newcommand{\nconn}{\mathrel{\not\leftrightarrow}}
\newcommand{\connin}[1]{\xleftrightarrow{\;#1\;}}
\newcommand{\ptwo}{p_{2\to 2}}
\newcommand{\pexp}{p_{\mathrm{exp}}}
\newcommand{\dbr}{d_{\mathrm{br}}}
\newcommand{\diambr}{\diam_{\mathrm{br}}}
\newcommand{\radbr}{\operatorname{rad}_{\mathrm{br}}}
\newcommand{\radint}{\operatorname{rad}_{\mathrm{int}}}
\newcommand{\radext}{\operatorname{rad}_{\mathrm{ext}}}
\newcommand{\Nb}{\mathrm{N}}                   % closed neighbourhood N[A]
\newcommand{\Cor}{\mathrm{Cor}}                % corridor event
\newcommand{\Bl}{\mathrm{Bl}}                  % blockers
\newcommand{\Erw}{\mathbf{E}^{\mathrm{rw}}_{k}}  % average over k-step lazy walk paths from o
\newcommand{\Ehat}{\widehat{\EE}}
